\section{三次縮約からの指数成長障害}
\label{sec:cubic-slice-growth}
有限次数のTorとは別に、商の経路成長から全ての候補に共通の排除条件を得る。
ここで$C$は四窓の元の二次代数とし、$E$はその二次双対である。

\begin{proposition}[一つの矢空間を一次元へ落とす成長検査]\label{prop:cut-arrow-growth}
切る位置を3とし、$0\ne\ell\in V_3^*$について
\[
 X=V_0/(\ell\otimes1)R_3,\qquad Z=V_2/(1\otimes\ell)R_2
\]
と置く。$R_0,R_1$の像を$\bar R_0\subset X\otimes V_1$,
$\bar R_1\subset V_1\otimes Z$とし、
\[
 d_\ell=\dim \frac{X\otimes V_1\otimes Z}
              {\bar R_0\otimes Z+X\otimes\bar R_1}
\]
とする。各$m\ge1$で$\dim C_{0,4m}\ge d_\ell^m$である。
特に標準AS対象では四位置の全ての$\ell$で$d_\ell\le1$である。
\end{proposition}
\begin{proof}
各周期の一次空間を順に$X,V_1,Z,k$へ写し、最後の商写像に$\ell$を使う。
二次関係の像で商を取れば$C$の商代数を得る。位置2,3の二次関係の像は零であり、
残る二つは一つの三文字ブロック内だけに作用する。従って長さ$4m$の経路商は
表示した三文字空間の$m$個のテンソル積である。
特に$d_\ell\ge2$では長さ64の下界$2^{16}=65536$が$\binom{67}{3}=47905$を超える。
標準AS対象の長さ$4m$の次元は$\binom{4m+3}{3}$なので結論を得る。
\end{proof}

\begin{theorem}[連続する二つの純粋関係の排除]\label{thm:no-consecutive-pure-relations}
標準有限Frobenius候補で、非零ベクトル$a,b,c$が
$a\otimes b\in R_0$, $b\otimes c\in R_1$を満たすなら、
$\dim C_{0,4m}\ge6^m$である。特に、標準AS対象のどの連続三窓でも
$W_{i,3}$は非零の可分解三テンソルを含まない。
\end{theorem}
\begin{proof}
$a\otimes b\otimes c$は三次交叉に属する。完全trace対合により
$V_3^*\to W_{0,3}$は同型なので、このテンソルはある$\ell$で$w$を縮約したものになる。
$R_3$を$\ell$で縮約した像はその三テンソルの先頭の支持$ka$、
$R_2$での像は末尾の支持$kc$である。これは$R_2,R_3$が最上テンソルの隣接縮約像だからである。
従って$\dim X=\dim Z=3$で、指定した純粋関係は各々$\bar R_0,\bar R_1$への射影で消える。
両像の次元は5以下であり、$d_\ell\ge36-5\cdot3-3\cdot5=6$となる。
逆に可分解三テンソルが交叉に入れば両端で縮約して指定の二関係を得る。
純粋関係のある始点では長さ16ですでに$6^4=1296>\binom{19}{3}=969$である。
どの位置でも命題\ref{prop:cut-arrow-growth}を適用できる。
\end{proof}

\begin{corollary}[正証明書のある対への階数1写像]\label{cor:no-rank-one-into-positive-pair}
標準AS対象で位置$i,j$を可逆な正証明書により同一視できるとする。
別の位置$b$から$i$への正symmetrizer空間は、階数1の写像を含まない。
ここで指定した二位置が隣接することは要らない。
\end{corollary}
\begin{proof}
$i,j$を同一視してその交換で$w$を不変とする。
階数1の$H=u\ell$が$b$から$i$への正証明書なら、
$H$を位置$b$に作用させたテンソルは位置$i,b$に対称であり、変換した側は$ku$に支持される。
従って$\ell$による元の位置$b$の縮約は位置$i$で$ku$に支持される。
$i,j$の対称性から位置$j$でも同じ直線に支持され、縮約は非零の可分解三テンソルとなる。
これは定理\ref{thm:no-consecutive-pure-relations}に反する。
\end{proof}

\begin{corollary}[三因子対称な基点枝の全次数成長]\label{cor:basepoint-exponential-growth}
三因子対称な標準有限Frobenius候補に基点$[p]$があるとき、
対称な連続三位置から始まる長さ$4m$の次元は少なくとも$9^m$である。
基点で二次式の射影Jacobian階数が3なら、下界は$10^m$になる。
Boolean境界$t=0$では、構成する商の次元はちょうど$10^m$である。
\end{corollary}
\begin{proof}
三次双対空間に$p^{\otimes3}$があり、対称な三因子の両端の商はいずれも$H=V/kp$である。
二つの内側の関係は同じ$R=B_2^*\subset\operatorname{Sym}^2 V$であり、
$\dim R=6$, $p^2\in R$である。片側の射影の核は$kp^2$である
（$p\otimes v$が対称なら$v\in kp$）。
$H\otimes V\otimes H$の両端反転を$\tau$とすると、挿入関係は
$J+\tau J$、$J=\bar R\otimes H$、$\dim J=15$である。
反転の正固有部分は周囲で$6\cdot4=24$次元、関係側で
$(1+\tau)J$の次元は15以下なので、商次元は9以上となる。

Jacobian階数3なら、補空間を選び$V=H\oplus kp$と書いたとき、
$R\to\operatorname{Sym}^2 H$の核も$kp^2$である。
実際、$p\odot v$の形の関係は全$q\in L$について$q(p,v)=0$を満たすので、
$v$の$H$成分が零になる。像を$K$とすると$\dim K=5$であり、
$\bar R\subset H\otimes V$は$K$から$H\otimes kp$への写像のグラフになる。
三文字の中央も$H$へ射影すると、
$J,\tau J$はそれぞれ$K\otimes H$, $H\otimes K$の上のグラフである。
両底空間の交叉は$\operatorname{Sym}^3 H$の中で、
$K^\perp$の生成二次式による縮約が零となる部分である。
その次元は$10-3=7$である。実際、双対写像は非零二次式による三変数一次式の乗法で単射である。
この交叉から二つのグラフへ持ち上げると、中央の$p$成分の差は両端の反転で符号が変わる。
その値域は$\bigwedge^2 H$の三次元空間なので、$\dim(J\cap\tau J)\ge7-3=4$である。
従って$d_\ell=36-30+\dim(J\cap\tau J)\ge10$である。
Boolean点の射影挿入行列では階数が26となることを有理数で直接検算できる。
ここで求めた等号は商に関するものである。
\end{proof}

\paragraph{射影的な有限検査。}
標準有限Frobenius候補の元代数は$\dim C_{0,2}=\dim C_{1,3}=10$、
$\dim C_{0,3}=64-24-24+4=20$を満たす。
命題\ref{prop:cut-arrow-growth}の$d_\ell$は次の写像の余核次元でもある。
\[
 R_3\otimes C_{1,3}\ \oplus\ C_{0,2}\otimes R_2
 \ \xrightarrow{\ \Phi(\ell)\ }\ C_{0,3},\qquad
 (r\otimes z,y\otimes r')\longmapsto
 ((\ell\otimes1)r)z+y((1\otimes\ell)r').
\]
実際、二つの内側の関係は$C_{0,3}$の段階ですでに割られ、残りの商は両端の線型関係による。
したがって$\Phi$は$20\times120$の一次行列である。ある位置で19次小行列式全てが
$\mathbb P(V_3^*)$に共通零点を持つ候補は指数成長を持つ。
この条件は添字保持同型で不変であり、候補族上で相対閉である。
実際、定数階数の経路商束の上の階数条件を射影束から射影すれば閉となり、位置は四つだけである。
これは非正則性の十分条件であり、全ての非Koszul候補の検出は示していない。
\path{computations/verify_cubic_slice_growth.py}に、Boolean点および正則対照点での商の検算を置く。

\begin{corollary}[対称閉包内の二対二の可逆証明書]\label{cor:two-pairs-transfer-ranks}
$Z_{\rm sym}$内の標準AS候補で、四位置が二つの可逆正証明書の対に分かれ、
異なる対の間には可逆な正証明書がないとする。異なる対から一つずつ位置を選ぶと、
その二位置間に両向きの正証明書$H,K$が存在し、
\[
 \rk H=\rk K=2,\qquad \operatorname{im}H=\ker K,\quad
 \operatorname{im}K=\ker H
\]
となる。
\end{corollary}
\begin{proof}
$S$上では全位置を対称化する二方向の同一視は射影的に互いに逆である。
それらの二つの射影行列を記録したグラフの閉包は、射影行列空間の積からの射影の固有性により
$Z_{\rm sym}$へ全射する。従って所与の点には、互換の正証明書方程式と
$HK,KH$がスカラーであるという閉条件を満たす非零の極限$H,K$が存在する。
対をまたぐため両者は特異で、その積のスカラーは零になる。
両方の標的位置がそれぞれ可逆対の中にあるので、系\ref{cor:no-rank-one-into-positive-pair}が
いずれの階数1も排除する。二階数の和は4以下であり、全ての等式を得る。
この結論は当該配置のAS点の存在を仮定する条件付きの制約である。
\end{proof}
